% modules/{{SLUG}}.tex — a manuscript module.
%
% Every claim-bearing theorem-environment \label here is exactly one ledger node whose
% `kind` is that environment; a \section or equation label is structural and is not a
% node. The ten claim environments are the ten `kind` values, declared in preamble.tex.
%
% A worked module holding an obstruction, an open question, a proved proposition and a
% refuted conjecture is example/modules/00-overview.tex.
\documentclass[../main.tex]{subfiles}
\begin{document}

\section{{{TITLE}}}
\label{sec:{{SLUG}}}

Orientation prose. What this module studies and which conventions everything in it
rests on. A fixed normalization — a sign, a scaling, a log base, which constant absorbs
what — is not prose: state it in a \texttt{definition} environment under its own
\texttt{\textbackslash label} and give it a ledger node with \texttt{kind: definition}
and \texttt{status: defined}, so that changing it is a mathematical edit the checker can
see.

\begin{{{KIND}}}
\label{{{NODE}}}
  The precise quantified statement. This text is canonical: the ledger's \texttt{summary}
  is a gloss of it, and the problem brief may quote it verbatim as a marked copy, but
  nothing else in the repository holds it.
\end{{{KIND}}}

\end{document}
